% ============================================================================
%  Laboratorio di Programmazione -- Lezione 4 (6 ottobre 2026)
%  Compilare con:  pdflatex -shell-escape Lezione04_Casting_Wrapper_Arrays_Modularita.tex
%  Richiede: mermaid.sty (nel progetto) + Mermaid CLI `mmdc` nel PATH
%  NB: ogni frame che contiene codice (lstlisting) o un diagramma (mermaid)
%      deve essere dichiarato [fragile].
% ============================================================================
\documentclass[10pt,aspectratio=169]{beamer}
\usetheme{Madrid}
\usecolortheme{default}

% `provide=*` carica la localizzazione italiana dai file .ini di babel: funziona anche
% su installazioni TeX prive di texlive-lang-italian; con una TeX Live completa e`
% equivalente al classico \usepackage[italian]{babel}.
\usepackage[italian,provide=*]{babel}
\usepackage[utf8]{inputenc}

% Localizzazione dei nomi dei blocchi Beamer
\uselanguage{italian}
\languagepath{italian}
\deftranslation[to=italian]{Definition}{Definizione}
\deftranslation[to=italian]{Theorem}{Teorema}
\deftranslation[to=italian]{Example}{Esempio}
\deftranslation[to=italian]{Proof}{Dimostrazione}
\deftranslation[to=italian]{Corollary}{Corollario}
\deftranslation[to=italian]{Lemma}{Lemma}

\usepackage{listings}
\usepackage{xcolor}
\usepackage{hyperref}
\usepackage{array}
\usepackage{booktabs}
\usepackage{mermaid}

% Mermaid: sfondo trasparente e dimensione massima adatta a una slide 16:9
\MermaidRendererOptions{-b transparent}
\MermaidAdjustBoxOpts{max width=0.95\linewidth,max height=0.62\textheight,center}

% Configurazione colori per il codice
\definecolor{codegreen}{rgb}{0,0.6,0}
\definecolor{codegray}{rgb}{0.5,0.5,0.5}
\definecolor{codepurple}{rgb}{0.58,0,0.82}
\definecolor{backcolour}{rgb}{0.95,0.95,0.92}
\definecolor{keywordblue}{rgb}{0.0,0.0,0.7}

% Stile per snippet Java
\lstdefinestyle{javastyle}{
    backgroundcolor=\color{backcolour},
    commentstyle=\color{codegreen}\itshape,
    keywordstyle=\color{keywordblue}\bfseries,
    numberstyle=\scriptsize\color{codegray},
    stringstyle=\color{codepurple},
    basicstyle=\ttfamily\footnotesize,
    breakatwhitespace=false,
    breaklines=true,
    captionpos=b,
    keepspaces=true,
    numbers=none,
    numbersep=5pt,
    showspaces=false,
    showstringspaces=false,
    showtabs=false,
    tabsize=2,
    language=Java,
    literate={à}{{\`a}}1 {è}{{\`e}}1 {é}{{\'e}}1 {ì}{{\`i}}1
             {ò}{{\`o}}1 {ù}{{\`u}}1 {È}{{\`E}}1
}
\lstset{style=javastyle}

% Stile per righe di terminale
\lstdefinestyle{shell}{
    language=bash,
    backgroundcolor=\color{black!85},
    basicstyle=\ttfamily\footnotesize\color{white},
    keywordstyle=\color{white},
    stringstyle=\color{white},
    commentstyle=\color{green!70},
    numbers=none,
    frame=none,
    showstringspaces=false
}

\title[Lezione 4 -- Casting, wrapper, Arrays, modularità]{Lezione 4 -- Casting, wrapper, \texttt{Arrays}, modularità}
\subtitle{Conversioni di tipo, classi wrapper, la classe \texttt{Arrays} e la programmazione modulare}
\author[Mattia Fonisto]{Mattia Fonisto}
\institute[Lab.\ di Programmazione]{Laboratorio di Programmazione -- Ingegneria Informatica, II anno\\
           Università degli Studi di Napoli Federico II -- Sede di San Giovanni a Teduccio}
\date{6 ottobre 2026}

\begin{document}

\frame{\titlepage}

% ========================================================================
\begin{frame}{Agenda della lezione (2h)}
\begin{enumerate}
    \item \textbf{Ripasso e Homework 3} \hfill (\emph{15 min})
    \begin{itemize}
        \item uno di voi alla lavagna con \texttt{Statistiche}, \texttt{Istogramma} o \texttt{Simmetrica}: lo commentiamo insieme
    \end{itemize}
    \item \textbf{Tipi: aspetti avanzati} \hfill (\emph{20 min})
    \begin{itemize}
        \item conversioni implicite e cast esplicito; \texttt{char} come numero (Unicode)
        \item classi wrapper, autoboxing e auto-unboxing; la classe \texttt{Character}
    \end{itemize}
    \item \textbf{\texttt{Arrays} e varargs} \hfill (\emph{10 min})
    \begin{itemize}
        \item stampare, copiare, confrontare, ordinare, cercare; argomenti a lunghezza variabile
    \end{itemize}
    \item \textbf{Programmazione modulare} \hfill (\emph{15 min})
    \begin{itemize}
        \item astrazione, modulo (interfaccia e corpo), information hiding, coesione e accoppiamento
    \end{itemize}
    \item \textbf{Esercitazione} -- tutti insieme \hfill (\emph{50 min})
    \begin{itemize}
        \item \emph{Commentiamo insieme}: \texttt{Statistiche}, un modulo con \texttt{Arrays} e varargs
        \item \emph{Insieme, passo passo}: \texttt{Frequenze} -- da un \texttt{main} che fa tutto a moduli coesi
    \end{itemize}
    \item \textbf{Wrap-up e Homework 4} \hfill (\emph{10 min})
\end{enumerate}
\end{frame}

% ========================================================================
\section{Casting e char}

\begin{frame}[fragile]{Conversioni implicite (\emph{widening}) -- \texttt{Widening.java}}
\begin{columns}[T]
\column{0.5\textwidth}
\begin{lstlisting}[basicstyle=\ttfamily\scriptsize]
public static double meta(double x) {
    return x / 2;
}

public static void main(String[] args) {
    char c = 'A';
    int i = c;       // char -> int: 65
    double d = i;    // int -> double: 65.0
    System.out.println(i + " " + d);

    System.out.println(c + 1);     // 66
    System.out.println(i + 0.5);   // 65.5
    System.out.println(7 / 2);     // 3
    System.out.println(7 / 2.0);   // 3.5
    System.out.println(meta(7));   // 3.5

    // int n = d;    // ERRORE: double -> int
    // char k = i;   // ERRORE: int -> char
}
\end{lstlisting}

\column{0.5\textwidth}
\MermaidAdjustBoxOpts{max width=\linewidth,max height=0.16\textheight,center}
\begin{mermaid}
flowchart LR
  c["<b>char</b><br/>16 bit"] ==>|"implicita"| i["<b>int</b><br/>32 bit"]
  i ==>|"implicita"| d["<b>double</b><br/>64 bit"]
  style c fill:#dcfce7,stroke:#15803d
  style i fill:#dbeafe,stroke:#1d4ed8
  style d fill:#fef3c7,stroke:#b45309
\end{mermaid}
\MermaidAdjustBoxOpts{max width=0.95\linewidth,max height=0.62\textheight,center}
\small
Verso un tipo che contiene \emph{tutti} i valori di quello di partenza la conversione è automatica:
\begin{itemize}\itemsep=2pt
    \item nell'\textbf{assegnazione}: \texttt{int i = c}, \texttt{double d = i}
    \item nelle \textbf{espressioni}: con operandi di tipo diverso vince il tipo più \textbf{capiente}. \texttt{char + int} è un \texttt{int}; \texttt{int + double} è un \texttt{double}
    \item nella \textbf{chiamata}: \texttt{meta(7)} riceve \texttt{7.0}
\end{itemize}
\vspace{0.1cm}
Nel verso opposto il compilatore si rifiuta. \texttt{boolean} non si converte in nessun altro tipo, in nessun modo.
\end{columns}
\end{frame}

% ========================================================================
\begin{frame}[fragile]{Il cast esplicito (\emph{narrowing}) -- \texttt{Narrowing.java}}
\begin{columns}[T]
\column{0.57\textwidth}
\begin{lstlisting}[basicstyle=\ttfamily\scriptsize]
double prezzo = 9.99;
int euro = (int) prezzo;            // 9: tronca
int vicino = (int) (prezzo + 0.5);  // 10
System.out.println(euro + " " + vicino);
System.out.println((int) -9.99);    // -9

int somma = 59;
int n = 6;
System.out.println(somma / n);             // 9
System.out.println((double) somma / n);    // 9.83...
System.out.println((double) (somma / n));  // 9.0

System.out.println((int) 3e10);     // 2147483647
char strano = (char) 2000000000;
System.out.println((int) strano);   // 37888
\end{lstlisting}
\small
\texttt{(tipo) espressione}: la conversione la chiedo io, e la perdita di informazione è una mia responsabilità.

\column{0.43\textwidth}
\small
\textbf{Da \texttt{double} a \texttt{int}} la parte decimale viene \textbf{troncata} (verso lo zero), non arrotondata. Per l'intero più vicino a un numero positivo: \texttt{(int) (x + 0.5)}.

\vspace{0.2cm}
\textbf{Il cast agisce su ciò che lo segue.} \texttt{(double) somma / n} converte \texttt{somma} e poi divide: divisione reale. \texttt{(double) (somma / n)} divide fra interi e converte il risultato: troppo tardi.

\vspace{0.2cm}
\textbf{Fuori dai limiti} nessun errore, nessun avviso:
\begin{itemize}\itemsep=2pt
    \item \texttt{double} \textrightarrow{} \texttt{int}: si ottiene il limite, \texttt{Integer.MAX\_VALUE}
    \item \texttt{int} \textrightarrow{} \texttt{char}: restano i 16 bit meno significativi
\end{itemize}
\end{columns}
\end{frame}

% ========================================================================
\begin{frame}[fragile]{\texttt{char} è un numero: il codice Unicode -- \texttt{Caratteri.java}}
\begin{columns}[T]
\column{0.54\textwidth}
\begin{lstlisting}[basicstyle=\ttfamily\scriptsize]
char c = 'a';
System.out.println((int) c);          // 97
System.out.println((char) 98);        // b
System.out.println(c + 1);            // 98: int
System.out.println((char) (c + 1));   // b

// c = c + 1;   // ERRORE: int -> char
c++;            // ok: ora c vale 'b'
c += 2;         // ok: ora c vale 'd'
System.out.println(c);                // d

char accentata = '\u00E8';   // la e accentata
System.out.println((int) accentata);  // 232

System.out.println('7' - '0');        // 7
System.out.println('e' - 'a');        // 4
System.out.println((char) ('a' + 4)); // e
System.out.println((char) ('e' - 'a' + 'A')); // E
\end{lstlisting}

\column{0.46\textwidth}
\small
Un \texttt{char} è un intero senza segno a 16 bit: il \textbf{codice Unicode} del carattere. \texttt{'a'} e \texttt{97} sono lo stesso valore.

\begin{center}\footnotesize
\begin{tabular}{lll}
\toprule
\texttt{'0'}\ldots\texttt{'9'} & \texttt{'A'}\ldots\texttt{'Z'} & \texttt{'a'}\ldots\texttt{'z'} \\
48\ldots 57 & 65\ldots 90 & 97\ldots 122 \\
\bottomrule
\end{tabular}
\end{center}

\begin{itemize}\itemsep=2pt
    \item cifre e lettere hanno codici \textbf{consecutivi}: \texttt{c - '0'} è il valore di una cifra, \texttt{c - 'a'} la posizione di una lettera, \texttt{(char) ('a' + k)} la lettera di posizione \texttt{k}
    \item \texttt{c + 1} è un \texttt{int}: per riavere un \texttt{char} serve il cast. \texttt{c++} e \texttt{c += 2} lo contengono già
    \item \texttt{'\textbackslash uXXXX'}: un carattere scritto con il suo codice, in esadecimale
\end{itemize}
\end{columns}
\end{frame}

% ========================================================================
\section{Classi wrapper}

\begin{frame}[fragile]{Le classi wrapper -- \texttt{Wrapper.java}}
\begin{columns}[T]
\column{0.55\textwidth}
\begin{lstlisting}[basicstyle=\ttfamily\scriptsize]
int n = 10;                       // il valore
Integer w = Integer.valueOf(10);  // un oggetto

// metodi di istanza: w e` un oggetto
System.out.println(w.intValue() + n);       // 20
System.out.println(w.doubleValue() / 4);    // 2.5
System.out.println(w.toString().length());  // 2

// costanti e metodi static: il corredo del tipo
int max = Integer.MAX_VALUE;   // 2147483647
int k = Integer.parseInt("42");           // 42
double x = Double.parseDouble("2.5");     // 2.5
String t = Integer.toString(42);          // "42"
String bin = Integer.toBinaryString(10);  // "1010"
System.out.println(max + " " + k + " " + x);
System.out.println(t + 1 + " " + bin);  // 421 1010
\end{lstlisting}
\small
Sono \textbf{immutabili}, come \texttt{String}: nessun metodo cambia il valore contenuto.

\column{0.45\textwidth}
\MermaidAdjustBoxOpts{max width=\linewidth,max height=0.2\textheight,center}
\begin{mermaid}
flowchart LR
  n["<b>stack</b><br/>int n<br/>10"]
  w["<b>stack</b><br/>Integer w<br/>●"]
  o["<b>heap</b><br/>oggetto Integer<br/>10"]
  n ~~~ w
  w --> o
  style n fill:#fef3c7,stroke:#b45309
  style w fill:#fef3c7,stroke:#b45309
  style o fill:#dbeafe,stroke:#1d4ed8
\end{mermaid}
\MermaidAdjustBoxOpts{max width=0.95\linewidth,max height=0.62\textheight,center}
\small
Ogni tipo primitivo ha la sua classe \emph{wrapper} (in \texttt{java.lang}): \texttt{Boolean}, \texttt{Character}, \texttt{Integer}, \texttt{Double}, e così per gli altri quattro.

\vspace{0.1cm}
\textbf{A cosa servono}
\begin{itemize}\itemsep=2pt
    \item a usare un valore primitivo dove sono ammessi solo \textbf{oggetti} (le strutture dati delle lezioni 11 e 12)
    \item i primitivi non hanno metodi: costanti e metodi \texttt{static} che li riguardano stanno nel wrapper
\end{itemize}
\end{columns}
\end{frame}

% ========================================================================
\begin{frame}[fragile]{Autoboxing e auto-unboxing -- \texttt{Boxing.java}}
\begin{columns}[T]
\column{0.52\textwidth}
\begin{lstlisting}[basicstyle=\ttfamily\scriptsize]
public static void azzera(Integer x) {
    x = 0;    // cambia solo la copia
}

public static void main(String[] args) {
    Integer w = 10;     // boxing
    int n = w;          // unboxing
    Integer s = w + n;  // unboxing e boxing
    System.out.println(s);             // 20

    azzera(w);
    System.out.println(w);             // 10

    Integer a = 1000;
    Integer b = 1000;
    System.out.println(a == b);        // false
    System.out.println(a.equals(b));   // true

    Integer[] voti = new Integer[3];  // 3 null!
    voti[0] = 28;                     // boxing
    int primo = voti[0];              // unboxing
    System.out.println(primo);        // 28
    // int x = voti[1];   // NullPointerException
}
\end{lstlisting}

\column{0.48\textwidth}
\MermaidAdjustBoxOpts{max width=\linewidth,max height=0.2\textheight,center}
\begin{mermaid}
flowchart LR
  v["voti ●"] --> arr["[0] ● &nbsp; [1] null &nbsp; [2] null"]
  arr --> o["Integer<br/>28"]
  style v fill:#fef3c7,stroke:#b45309
  style arr fill:#dbeafe,stroke:#1d4ed8
  style o fill:#dcfce7,stroke:#15803d
\end{mermaid}
\MermaidAdjustBoxOpts{max width=0.95\linewidth,max height=0.62\textheight,center}
\footnotesize
Il compilatore converte da solo fra primitivo e wrapper: \textbf{boxing} (\texttt{Integer.valueOf(10)}) e \textbf{unboxing} (\texttt{w.intValue()}).
\begin{itemize}\itemsep=2pt
    \item restano \textbf{oggetti}: \texttt{==} confronta i riferimenti, il valore si confronta con \texttt{equals}. Fra $-128$ e $127$ \texttt{==} fra \texttt{Integer} dà \texttt{true} (una cache, come il pool delle \texttt{String}); fra \texttt{Double} mai. Fra primitivi \texttt{==} resta corretto
    \item un \texttt{Integer[]} nasce pieno di \texttt{null}, non di zeri (un \texttt{int[]} sì): l'unboxing di \texttt{null} ferma il programma
    \item \texttt{azzera} non cambia \texttt{w}: riceve una copia del riferimento, e \texttt{x = 0} la fa puntare a un altro oggetto
\end{itemize}
\end{columns}
\end{frame}

% ========================================================================
\begin{frame}[fragile]{La classe \texttt{Character} -- \texttt{InfoCarattere.java}}
\begin{columns}[T]
\column{0.54\textwidth}
\begin{lstlisting}[basicstyle=\ttfamily\scriptsize]
String s = "Aula T2, ore 9:30";
int lettere = 0;
int cifre = 0;
int maiuscole = 0;
for (int i = 0; i < s.length(); i++) {
    char c = s.charAt(i);
    if (Character.isLetter(c)) {
        lettere++;
    }
    if (Character.isDigit(c)) {
        cifre++;
    }
    if (Character.isUpperCase(c)) {
        maiuscole++;
    }
}
System.out.println(lettere + " lettere, " + cifre
        + " cifre, " + maiuscole + " maiuscole");

char m = Character.toUpperCase('a');       // 'A'
char r = Character.toUpperCase('7');       // '7'
System.out.println(m + " " + r);
\end{lstlisting}

\column{0.46\textwidth}
\begin{lstlisting}[style=shell,basicstyle=\ttfamily\scriptsize\color{white}]
8 lettere, 4 cifre, 2 maiuscole
A 7
\end{lstlisting}
\small
Il wrapper di \texttt{char}. Quasi tutti i suoi metodi sono \texttt{static} e ricevono il carattere come argomento:

\begin{center}\scriptsize
\begin{tabular}{@{}ll@{}}
\toprule
lettera? cifra? & \texttt{isLetter}, \texttt{isDigit} \\
lettera o cifra? & \texttt{isLetterOrDigit} \\
maiuscola? minuscola? & \texttt{isUpperCase}, \texttt{isLowerCase} \\
spazio, tabulazione, a capo? & \texttt{isWhitespace} \\
conversione & \texttt{toUpperCase}, \texttt{toLowerCase} \\
\bottomrule
\end{tabular}
\end{center}

\begin{itemize}\itemsep=2pt
    \item \texttt{toUpperCase} restituisce l'argomento stesso se non esiste una maiuscola corrispondente
    \item \texttt{isLetter} conosce \emph{tutte} le lettere Unicode, accentate comprese; \texttt{c >= 'a' \&\& c <= 'z'} solo quelle 26
\end{itemize}
\end{columns}
\end{frame}

% ========================================================================
\section{Arrays e varargs}

\begin{frame}[fragile]{\texttt{java.util.Arrays}: stampare, copiare, confrontare -- \texttt{UsaArrays.java}}
\begin{columns}[T]
\column{0.56\textwidth}
\begin{lstlisting}[basicstyle=\ttfamily\scriptsize]
import java.util.Arrays;

int[] a = {4, 5, 1, 13, 0};
System.out.println(a);            // riferimento!
System.out.println(Arrays.toString(a));

int[] b = Arrays.copyOf(a, a.length);  // copia vera
b[0] = 99;
System.out.println(a[0]);         // a non cambia
int[] lungo = Arrays.copyOf(a, 7);
System.out.println(Arrays.toString(lungo));

int[] c = {4, 5, 1, 13, 0};
System.out.println(a == c);
System.out.println(a.equals(c));  // come ==
System.out.println(Arrays.equals(a, c));

int[] d = new int[4];
Arrays.fill(d, 7);
System.out.println(Arrays.toString(d));
\end{lstlisting}
\small
\texttt{Arrays} \textbf{non è un wrapper}: è una classe di soli metodi \texttt{static}, che lavorano sugli array ricevuti come argomento.

\column{0.44\textwidth}
\begin{lstlisting}[style=shell,basicstyle=\ttfamily\scriptsize\color{white}]
[I@659e0bfd
[4, 5, 1, 13, 0]
4
[4, 5, 1, 13, 0, 0, 0]
false
false
true
[7, 7, 7, 7]
\end{lstlisting}
\footnotesize
\begin{itemize}\itemsep=1pt
    \item \texttt{toString}: il contenuto come stringa
    \item \texttt{copyOf(a, n)}: un array \emph{nuovo} di \texttt{n} elementi (zeri in coda se \texttt{n} è maggiore)
    \item \texttt{Arrays.equals(a, c)}: confronta il contenuto
    \item \texttt{fill}: lo stesso valore in ogni elemento
\end{itemize}
\vspace{0.05cm}
\textbf{Perché \texttt{a.equals(c)} dà \texttt{false}?} \texttt{equals} confronta il contenuto solo nelle classi che lo definiscono così, come \texttt{String} e i wrapper. Gli array non lo fanno: il loro \texttt{equals} si comporta come \texttt{==}. Per questo esiste \texttt{Arrays.equals}.
\end{columns}
\end{frame}

% ========================================================================
\begin{frame}[fragile]{\texttt{Arrays}: ordinare e cercare -- \texttt{OrdinaCerca.java}}
\begin{columns}[T]
\column{0.52\textwidth}
\begin{lstlisting}[basicstyle=\ttfamily\scriptsize]
int[] v = {12, 7, 3, 21, 9};
Arrays.sort(v);        // MODIFICA v
System.out.println(Arrays.toString(v));
System.out.println(Arrays.binarySearch(v, 12));
System.out.println(Arrays.binarySearch(v, 8));

String[] nomi = {"Marta", "anna", "Luca"};
Arrays.sort(nomi);     // usa compareTo
System.out.println(Arrays.toString(nomi));

int[][] m = {{1, 2}, {3, 4}};
System.out.println(Arrays.toString(m));
System.out.println(Arrays.deepToString(m));
\end{lstlisting}
\begin{lstlisting}[style=shell,basicstyle=\ttfamily\scriptsize\color{white}]
[3, 7, 9, 12, 21]
3
-3
[Luca, Marta, anna]
[[I@15db9742, [I@6d06d69c]
[[1, 2], [3, 4]]
\end{lstlisting}

\column{0.48\textwidth}
\footnotesize
\begin{itemize}\itemsep=3pt
    \item \texttt{sort} ordina in senso crescente, \textbf{sul posto}: è \texttt{void} e modifica l'array del chiamante, attraverso la copia del riferimento (come \texttt{inverti} nella lezione~3). Con \texttt{String} e i wrapper, immutabili, nessun metodo può farlo
    \item per conservare l'ordine originale: \texttt{copyOf}, poi \texttt{sort} sulla copia
    \item \texttt{binarySearch} vuole un array \textbf{ordinato}: come cercare una parola nel dizionario, si apre a metà e si prosegue nella metà giusta. Restituisce l'indice, o un numero negativo se l'elemento non c'è (lezione~11)
    \item per le matrici: \texttt{deepToString} e \texttt{deepEquals}
\end{itemize}

\begin{center}\footnotesize
\begin{tabular}{ll}
\toprule
\emph{leggono} & \texttt{toString}, \texttt{equals}, \texttt{binarySearch} \\
\emph{modificano} & \texttt{sort}, \texttt{fill} \\
\emph{creano} & \texttt{copyOf} \\
\bottomrule
\end{tabular}
\end{center}
\end{columns}
\end{frame}

% ========================================================================
\begin{frame}[fragile]{Argomenti a lunghezza variabile (\emph{varargs}) -- \texttt{Varargs.java}}
\begin{columns}[T]
\column{0.53\textwidth}
\begin{lstlisting}[basicstyle=\ttfamily\scriptsize]
/** Pre: almeno un numero. */
public static double media(double... numeri) {
    double somma = 0;
    for (double x : numeri) {   // un double[]
        somma += x;
    }
    return somma / numeri.length;
}

/** Il primo argomento (a) e` obbligatorio. */
public static int massimo(int a, int... altri) {
    int max = a;
    for (int x : altri) {
        if (x > max) {
            max = x;
        }
    }
    return max;
}
\end{lstlisting}
\small
\texttt{tipo...\ nome}: il metodo accetta \textbf{zero o più} argomenti di quel tipo. Dentro il metodo \texttt{nome} è un normale array.

\column{0.47\textwidth}
\begin{lstlisting}[basicstyle=\ttfamily\scriptsize]
double[] voti = {28, 30, 24};
double a = media(10, 40);          // 25.0
double b = media(10, 40, 50, 75);  // 43.75
double c = media(voti);            // 27.33
double d = media();                // NaN
int m = massimo(3, 9, 4);          // 9
// massimo();    // ERRORE: manca il primo
System.out.println(a + " " + b + " " + c);
System.out.println(d + " " + m);
\end{lstlisting}
\footnotesize
\begin{itemize}\itemsep=2pt
    \item variabile è il \textbf{numero di argomenti}, non l'array: a ogni chiamata il compilatore ne crea uno della lunghezza giusta, che poi, come ogni array, non cambia più. Se ne può anche passare uno già pronto
    \item un solo parametro varargs, e \textbf{per ultimo}
    \item \texttt{media()} compila, e \texttt{0.0 / 0} dà \texttt{NaN}: un parametro normale davanti, come in \texttt{massimo}, rende obbligatorio il primo argomento
    \item \texttt{printf} è un metodo varargs
\end{itemize}
\end{columns}
\end{frame}

% ========================================================================
\section{Programmazione modulare}

\begin{frame}[fragile]{Programmazione modulare: suddividere un sistema in parti}
\begin{columns}[T]
\column{0.57\textwidth}
Un sistema software complesso non si scrive, né si capisce, tutto insieme: lo si \textbf{suddivide} in parti più piccole, i \emph{moduli}, ciascuna con un compito ben definito, e si stabilisce come queste parti sono collegate fra loro. L'insieme delle parti e dei loro collegamenti è l'\emph{architettura} del sistema.

\vspace{0.15cm}
\small
Ne usate dalla prima lezione: \texttt{Scanner}, \texttt{Math}, \texttt{String} sono moduli scritti da altri. Sapete cosa fanno e come si chiamano i loro metodi; il loro codice non l'avete mai letto.

\vspace{0.15cm}
\textbf{Che cosa si guadagna}
\begin{itemize}\itemsep=2pt
    \item \textbf{riuso}: lo stesso modulo serve a più programmi
    \item \textbf{modificabilità}: si corregge o si migliora un modulo senza toccare gli altri
    \item \textbf{leggibilità}: si capisce, e si prova, un pezzo alla volta
    \item \textbf{sviluppo separato}: più persone lavorano insieme, ciascuna sul suo modulo
\end{itemize}

\column{0.43\textwidth}
\MermaidAdjustBoxOpts{max width=\linewidth,max height=0.36\textheight,center}
\begin{mermaid}
flowchart TB
  P["<b>Playlist</b>"] ---|"usa"| R["<b>Riproduttore audio</b>"]
  P ---|"è composta da"| B["<b>Brano</b>"]
  style P fill:#dbeafe,stroke:#1d4ed8
  style R fill:#dcfce7,stroke:#15803d
  style B fill:#fef3c7,stroke:#b45309
\end{mermaid}
\MermaidAdjustBoxOpts{max width=0.95\linewidth,max height=0.62\textheight,center}
\small
Un lettore musicale suddiviso in tre moduli, e i due modi in cui due moduli possono essere collegati:
\begin{itemize}\itemsep=2pt
    \item \textbf{usa}: la playlist si serve del riproduttore audio per suonare
    \item \textbf{è composta da}: la playlist è fatta di brani
\end{itemize}
\end{columns}
\end{frame}

% ========================================================================
\begin{frame}[fragile]{Programmazione modulare: astrarre}
\begin{columns}[T]
\column{0.6\textwidth}
\begin{block}{Astrazione}
Considerare le proprietà \emph{rilevanti} di un'entità, ignorando i dettagli non essenziali.
\end{block}

\small
La stessa persona ha una vista \emph{anagrafica} (nome, data di nascita, residenza) e una vista \emph{clinica} (peso, pressione, temperatura). Nessuna delle due è ``tutta'' la persona: ciascuna tiene ciò che serve a chi la usa.
\begin{itemize}\itemsep=3pt
    \item \textbf{astrazione sul controllo}: una funzionalità, separata dai dettagli di come è realizzata. Sono i metodi: di \texttt{Arrays.sort} sapete \emph{cosa} fa, non \emph{come}
    \item \textbf{astrazione sui dati}: un dato, separato dai dettagli di come è memorizzato. Di una \texttt{String} sapete \emph{cosa} rappresenta (una sequenza di caratteri) e quali operazioni permette (\texttt{length}, \texttt{charAt}\ldots), non \emph{come} i caratteri sono tenuti in memoria. Sono i tipi: i nostri li definiremo con le classi
\end{itemize}

\column{0.4\textwidth}
\MermaidAdjustBoxOpts{max width=\linewidth,max height=0.42\textheight,center}
\begin{mermaid}
flowchart TB
  P["<b>una persona</b>"] -- "per l'anagrafe" --> A["<b>vista anagrafica</b><br/>nome<br/>data di nascita<br/>residenza"]
  P -- "per il medico" --> C["<b>vista clinica</b><br/>peso<br/>pressione<br/>temperatura"]
  style P fill:#fef3c7,stroke:#b45309
  style A fill:#dbeafe,stroke:#1d4ed8
  style C fill:#dcfce7,stroke:#15803d
\end{mermaid}
\MermaidAdjustBoxOpts{max width=0.95\linewidth,max height=0.62\textheight,center}
\small
L'astrazione fondamentale dei linguaggi procedurali, come il C, è quella sul controllo; quella dei linguaggi a oggetti, come Java, è l'astrazione sui dati.
\end{columns}
\end{frame}

% ========================================================================
\begin{frame}[fragile]{Programmazione modulare: il modulo}
\begin{columns}[T]
\column{0.64\textwidth}
\begin{block}{Modulo}
Una delle parti in cui il sistema è suddiviso: realizza un'astrazione. Chi lo usa, un altro modulo o un programma, è un suo \textbf{client}.
\end{block}

\small
L'astrazione può essere solo sul controllo (\texttt{Math}: operazioni) oppure sui dati e sul controllo insieme (\texttt{String}: un dato con le sue operazioni).

\vspace{0.08cm}
Un modulo \textbf{esporta} metodi e costanti, cioè li mette a disposizione dei client, e può \textbf{importare} quelli di altri moduli. È composto da due parti, tenute separate:
\begin{itemize}\itemsep=2pt
    \item l'\textbf{interfaccia}: \emph{cosa} fa il modulo e \emph{come si usa}. È la sola parte visibile ai client. Di \texttt{Math.sqrt} è \texttt{double sqrt(double a)}
    \item il \textbf{corpo}: \emph{come} l'astrazione è realizzata. Nascosto e protetto. Di \texttt{Math.sqrt} sono le istruzioni che calcolano la radice
\end{itemize}

\textbf{Information hiding}: al corpo si accede soltanto attraverso l'interfaccia, che espone lo stretto necessario e resta \textbf{stabile}. Finché non cambia, il corpo si può riscrivere senza toccare i client.

\column{0.36\textwidth}
\MermaidAdjustBoxOpts{max width=\linewidth,max height=0.57\textheight,center}
\begin{mermaid}
flowchart TB
  C["<b>Client</b><br/>chi usa il modulo"] -- "usa" --> I
  subgraph M[" "]
    I["<b>Interfaccia del modulo</b><br/>cosa fa, come si usa<br/><i>visibile ai client</i>"]
    B["<b>Corpo del modulo</b><br/>come è realizzato<br/><i>nascosto e protetto</i>"]
    I --- B
  end
  style C fill:#fef3c7,stroke:#b45309
  style I fill:#dcfce7,stroke:#15803d
  style B fill:#e5e7eb,stroke:#4b5563
\end{mermaid}
\MermaidAdjustBoxOpts{max width=0.95\linewidth,max height=0.62\textheight,center}
\footnotesize
Java ha anche una parola chiave \texttt{interface}: un costrutto specifico (lezione~9).
\end{columns}
\end{frame}

% ========================================================================
\begin{frame}[fragile]{Un modulo in Java, oggi -- \texttt{Voti.java} e \texttt{ProvaVoti.java}}
\begin{columns}[T]
\column{0.53\textwidth}
\begin{lstlisting}[basicstyle=\ttfamily\scriptsize]
/** Modulo: i voti d'esame in trentesimi. */
public class Voti {

    public static final int MINIMO = 18;
    public static final int MASSIMO = 30;

    /** true se v e` un voto valido. */
    public static boolean isValido(int v) {
        return v >= MINIMO && v <= MASSIMO;
    }

    /** La media riportata in 110-esimi. */
    public static int baseLaurea(double media) {
        return arrotonda(media * 110 / MASSIMO);
    }

    // corpo: visibile solo dentro Voti
    private static int arrotonda(double x) {
        return (int) (x + 0.5);
    }
}
\end{lstlisting}

\column{0.47\textwidth}
{\footnotesize Nel \texttt{main} di \texttt{ProvaVoti}:}
\begin{lstlisting}[basicstyle=\ttfamily\scriptsize]
boolean ok = Voti.isValido(31);    // false
int base = Voti.baseLaurea(27.4);  // 100
System.out.println(ok + " " + base);
System.out.println(Voti.MASSIMO);  // 30
// Voti.arrotonda(2.5);  // ERRORE: private
\end{lstlisting}
\begin{lstlisting}[style=shell,basicstyle=\ttfamily\scriptsize\color{white}]
$ javac ProvaVoti.java
$ java ProvaVoti
false 100
30
\end{lstlisting}
\footnotesize
\begin{itemize}\itemsep=2pt
    \item un modulo: una \textbf{classe} di metodi \texttt{static} nel suo file, \textbf{senza \texttt{main}}
    \item \texttt{public} è l'interfaccia: di ogni metodo al client bastano la firma e il commento \texttt{/** */}
    \item \textbf{\texttt{private}} è il corpo: si usa solo dentro la classe
    \item il client, qui \texttt{ProvaVoti}, scrive \texttt{Voti.metodo(...)}
    \item i due file stanno nella stessa cartella: compilando \texttt{ProvaVoti.java}, \texttt{javac} compila anche \texttt{Voti.java}
\end{itemize}
\end{columns}
\end{frame}

% ========================================================================
\begin{frame}[fragile]{Programmazione modulare: coesione e accoppiamento}
\begin{columns}[T]
\column{0.47\textwidth}
\begin{block}{Coesione (dentro un modulo)}
\small Un modulo è \textbf{fortemente coeso} se racchiude caratteristiche omogenee: realizza \emph{una sola} astrazione. \texttt{Voti} sa tutto dei voti e nient'altro; \texttt{Math} solo di matematica.
\end{block}
\column{0.47\textwidth}
\begin{block}{Accoppiamento (fra moduli)}
\small Misura quanto un modulo dipende dagli altri. Due moduli sono \textbf{debolmente accoppiati} se ciascuno conosce dell'altro solo l'interfaccia: gli passa dei parametri e ne riceve un risultato. Sono fortemente accoppiati se lavorano sugli stessi dati: non si può cambiarne uno senza toccare l'altro.
\end{block}
\end{columns}

\vspace{0.2cm}
\MermaidAdjustBoxOpts{max width=\linewidth,max height=0.2\textheight,center}
\begin{columns}[T]
\column{0.47\textwidth}
\centering\small accoppiamento \textbf{forte}
\begin{mermaid}
flowchart LR
  A["Modulo A"] <--> G["dati<br/>condivisi"]
  G <--> B["Modulo B"]
  style G fill:#fee2e2,stroke:#b91c1c
  style A fill:#dbeafe,stroke:#1d4ed8
  style B fill:#dbeafe,stroke:#1d4ed8
\end{mermaid}
\column{0.47\textwidth}
\centering\small accoppiamento \textbf{debole}
\begin{mermaid}
flowchart LR
  A["Modulo A"] -- "parametri" --> B["Modulo B"]
  B -. "risultato" .-> A
  style A fill:#dbeafe,stroke:#1d4ed8
  style B fill:#dbeafe,stroke:#1d4ed8
\end{mermaid}
\end{columns}
\MermaidAdjustBoxOpts{max width=0.95\linewidth,max height=0.62\textheight,center}

\vspace{0.15cm}
\small L'obiettivo è \textbf{alta coesione e basso accoppiamento}, ma i due criteri tirano in direzioni opposte: moduli piccolissimi sono molto coesi, ma per fare qualcosa devono usarsi a vicenda di continuo. Le dipendenze si moltiplicano e l'accoppiamento cresce. Va cercato un compromesso.
\end{frame}

% ========================================================================
\section{Esercitazione}

\begin{frame}[fragile]{Obiettivo dell'esercitazione}
Due programmi, tutti insieme, sulla stessa idea: \textbf{ogni file è un modulo con un compito solo}, e chi lo usa ne conosce soltanto l'interfaccia.

\MermaidAdjustBoxOpts{max width=0.95\linewidth,max height=0.4\textheight,center}
\begin{mermaid}
flowchart LR
  PS["<b>ProvaStatistiche</b><br/>main"] -- "usa" --> S["<b>Statistiche</b><br/>somma · media · minimo<br/>massimo · mediana"]
  F["<b>Frequenze</b><br/>main: solo ingresso e uscita"] -- "usa" --> S
  F -- "usa" --> L["<b>Lettere</b><br/>conta · lettera"]
  S -- "usa" --> A["<b>java.util.Arrays</b><br/>copyOf · sort"]
  style PS fill:#fef3c7,stroke:#b45309
  style F fill:#fef3c7,stroke:#b45309
  style S fill:#dcfce7,stroke:#15803d
  style L fill:#dcfce7,stroke:#15803d
  style A fill:#e5e7eb,stroke:#4b5563
\end{mermaid}
\MermaidAdjustBoxOpts{max width=0.95\linewidth,max height=0.62\textheight,center}

\begin{enumerate}\itemsep=4pt
    \item \textbf{Commentiamo insieme} (20') -- \texttt{Statistiche.java} e \texttt{ProvaStatistiche.java}: un modulo di statistiche su interi, con varargs, \texttt{Arrays} e cast. Lo leggiamo distinguendo interfaccia e corpo.
    \item \textbf{Insieme, passo passo} (30') -- \texttt{Frequenze}: partiamo da un \texttt{main} che fa tutto da solo e lo riscriviamo a moduli, riusando \texttt{Statistiche}. Io al PC proiettato, voi in VS Code.
\end{enumerate}
\end{frame}

% ========================================================================
\begin{frame}[fragile]{Commentiamo insieme -- \texttt{Statistiche.java} (1/2)}
\begin{columns}[T]
\column{0.5\textwidth}
\begin{lstlisting}[basicstyle=\ttfamily\scriptsize]
import java.util.Arrays;

/**
 * Modulo: statistiche su sequenze di interi.
 * Tutti i metodi LEGGONO gli argomenti.
 * Precondizione comune: almeno un valore.
 */
public class Statistiche {

    public static int somma(int... v) {
        int s = 0;
        for (int x : v) {
            s += x;
        }
        return s;
    }

    public static double media(int... v) {
        return (double) somma(v) / v.length;
    }
\end{lstlisting}

\column{0.5\textwidth}
\begin{lstlisting}[basicstyle=\ttfamily\scriptsize]
    public static int minimo(int... v) {
        int min = v[0];
        for (int x : v) {
            if (x < min) {
                min = x;
            }
        }
        return min;
    }

    public static int massimo(int... v) {
        int max = v[0];
        for (int x : v) {
            if (x > max) {
                max = x;
            }
        }
        return max;
    }
\end{lstlisting}
\end{columns}
\end{frame}

% ========================================================================
\begin{frame}[fragile]{Commentiamo insieme -- \texttt{Statistiche.java} (2/2)}
\begin{columns}[T]
\column{0.54\textwidth}
\begin{lstlisting}[basicstyle=\ttfamily\scriptsize]
    /** Il valore centrale dei dati ordinati. */
    public static double mediana(int... v) {
        int[] o = ordinata(v);
        int meta = o.length / 2;
        if (o.length % 2 == 1) {
            return o[meta];   // int -> double
        }
        return (o[meta - 1] + o[meta]) / 2.0;
    }

    // corpo: CREA una copia ordinata di v
    private static int[] ordinata(int[] v) {
        int[] copia = Arrays.copyOf(v, v.length);
        Arrays.sort(copia);
        return copia;
    }
}
\end{lstlisting}

\column{0.46\textwidth}
\small
\textbf{Da notare, leggendo:}
\begin{itemize}\itemsep=3pt
    \item \texttt{int...\ v}: ogni metodo si chiama con un array oppure con un elenco di valori
    \item \texttt{media}: il cast è su \texttt{somma(v)}, prima della divisione
    \item \texttt{mediana}: \texttt{return o[meta]} restituisce un \texttt{int} dove è atteso un \texttt{double} (conversione implicita); con un numero pari di valori, la media dei due centrali: \texttt{/ 2.0} e non \texttt{/ 2}
    \item \texttt{ordinata} è \texttt{private}: appartiene al corpo. \texttt{Arrays.sort} \emph{modifica}, il modulo promette di \emph{leggere}: si ordina una copia
    \item l'interfaccia: cinque metodi \texttt{public} e il commento iniziale
\end{itemize}
\end{columns}
\end{frame}

% ========================================================================
\begin{frame}[fragile]{Commentiamo insieme -- \texttt{ProvaStatistiche.java}}
\begin{columns}[T]
\column{0.55\textwidth}
\begin{lstlisting}[basicstyle=\ttfamily\scriptsize]
int[] voti = {28, 30, 24, 27, 30, 18};

System.out.println(Statistiche.somma(voti));
System.out.println(Statistiche.media(voti));
System.out.println(Statistiche.minimo(voti));
System.out.println(Statistiche.massimo(voti));
System.out.println(Statistiche.mediana(voti));
System.out.println(Arrays.toString(voti));

System.out.println(Statistiche.media(18, 30));
System.out.println(Statistiche.mediana(5, 1, 3));
int vicina = (int) (Statistiche.media(voti) + 0.5);
System.out.println(vicina);
\end{lstlisting}
\begin{lstlisting}[style=shell,basicstyle=\ttfamily\scriptsize\color{white}]
157
26.166666666666668
18
30
27.5
[28, 30, 24, 27, 30, 18]
24.0
3.0
26
\end{lstlisting}

\column{0.45\textwidth}
\small
\textbf{Rispondiamo insieme:}
\begin{enumerate}\itemsep=4pt
    \item \texttt{Statistiche.media(18, 30)}: chi costruisce l'array? Che tipo ha \texttt{v} dentro \texttt{media}?
    \item dopo \texttt{mediana(voti)} l'array \texttt{voti} è ancora nell'ordine originale: quale riga del modulo lo garantisce?
    \item se in \texttt{media} togliamo \texttt{(double)}, cosa stampa la seconda riga?
    \item possiamo scrivere qui \texttt{Statistiche.ordinata(voti)}? Perché?
    \item \texttt{Statistiche.media()} compila? E \texttt{Statistiche.minimo()}? Cosa succede eseguendoli?
\end{enumerate}
\end{columns}
\end{frame}

% ========================================================================
\begin{frame}[fragile]{Insieme, passo passo (30 min) -- \texttt{FrequenzeMonolite.java}}
\small
\textbf{Il problema.} Contare quante volte compare ogni lettera di un testo, in percentuale, e segnalare la più frequente.

\begin{columns}[T]
\column{0.52\textwidth}
\begin{lstlisting}[basicstyle=\ttfamily\scriptsize]
Scanner in = new Scanner(System.in);
System.out.print("Testo: ");
String testo = in.nextLine().toLowerCase();
in.close();
String alfabeto = "abcdefghijklmnopqrstuvwxyz";
int[] conteggi = new int[26];
int totale = 0;
for (int i = 0; i < testo.length(); i++) {
    int p = alfabeto.indexOf(testo.charAt(i));
    if (p >= 0) {
        conteggi[p]++;
        totale++;
    }
}
\end{lstlisting}
\small
Questo \texttt{main} funziona, ed è scritto con i soli strumenti delle lezioni 1--3. Ma fa \emph{tutto}: legge, conta le lettere, cerca il massimo, calcola le percentuali, stampa. Nessuna di queste parti si può riusare altrove, né provare da sola.

\column{0.48\textwidth}
\begin{lstlisting}[basicstyle=\ttfamily\scriptsize]
int max = conteggi[0];
for (int i = 1; i < 26; i++) {
    if (conteggi[i] > max) {
        max = conteggi[i];
    }
}
for (int i = 0; i < 26; i++) {
    int n = conteggi[i];
    if (n > 0) {
        double perc = 100.0 * n / totale;
        System.out.printf("%c %3d %5.1f%%",
            alfabeto.charAt(i), n, perc);
        if (n == max) {
            System.out.print("  <-- max");
        }
        System.out.println();
    }
}
\end{lstlisting}
\end{columns}
\end{frame}

% ========================================================================
\begin{frame}[fragile]{Insieme, passo passo -- dal \texttt{main} che fa tutto ai moduli}
\small
\textbf{Prima di scrivere codice, decidiamo insieme:}
\begin{enumerate}\itemsep=3pt
    \item quanti \textbf{compiti} diversi ha quel \texttt{main}? \textrightarrow{} ingresso e uscita; contare le lettere; sommare e trovare il massimo
    \item quali \textbf{moduli}? \textrightarrow{} \texttt{Lettere} (nuovo): trasforma le lettere in indici e le conta; \texttt{Statistiche}: è già pronto, lo \textbf{riusiamo}; \texttt{Frequenze}: il \texttt{main}, solo ingresso e uscita
    \item l'\textbf{interfaccia} di \texttt{Lettere}: cosa serve al client? \textrightarrow{} \texttt{int[] conta(String s)} e \texttt{char lettera(int indice)}; il resto è corpo, \texttt{private}
    \item cosa cambia con gli strumenti di oggi? \textrightarrow{} \texttt{c - 'a'} al posto di \texttt{alfabeto.indexOf(c)}; \texttt{(char) ('a' + i)} al posto di \texttt{alfabeto.charAt(i)}; il cast al posto di \texttt{100.0 *}
\end{enumerate}

\MermaidAdjustBoxOpts{max width=0.95\linewidth,max height=0.2\textheight,center}
\begin{mermaid}
flowchart LR
  c["char c = 'e'"] -->|"c - 'a'"| i["indice 4"]
  i -->|"conteggi[4]++"| a["int[26] conteggi"]
  a -->|"(char) ('a' + i)"| r["la lettera<br/>da stampare"]
  style c fill:#dcfce7,stroke:#15803d
  style r fill:#dcfce7,stroke:#15803d
  style i fill:#fef3c7,stroke:#b45309
  style a fill:#dbeafe,stroke:#1d4ed8
\end{mermaid}
\MermaidAdjustBoxOpts{max width=0.95\linewidth,max height=0.62\textheight,center}

Da tenere a mente: il programma nuovo deve stampare \emph{esattamente} ciò che stampa il vecchio.
\end{frame}

% ========================================================================
\begin{frame}[fragile]{Una possibile soluzione (1/2) -- il modulo \texttt{Lettere.java}}
\begin{columns}[T]
\column{0.5\textwidth}
\begin{lstlisting}[basicstyle=\ttfamily\scriptsize]
/** Modulo: le 26 lettere come indici. */
public class Lettere {

    public static final int QUANTE = 26;

    /** L'indice 0 e` 'a', il 25 e` 'z'. */
    public static char lettera(int indice) {
        return (char) ('a' + indice);
    }

    /** CREA l'array dei 26 conteggi. */
    public static int[] conta(String s) {
        int[] conteggi = new int[QUANTE];
        for (int i = 0; i < s.length(); i++) {
            char c = s.charAt(i);
            c = Character.toLowerCase(c);
            if (isLettera(c)) {
                conteggi[indice(c)]++;
            }
        }
        return conteggi;
    }
\end{lstlisting}

\column{0.5\textwidth}
\begin{lstlisting}[basicstyle=\ttfamily\scriptsize]
    // corpo: solo le minuscole da 'a' a 'z'
    private static boolean isLettera(char c) {
        return c >= 'a' && c <= 'z';
    }

    // 'a' -> 0, 'b' -> 1, ..., 'z' -> 25
    private static int indice(char c) {
        return c - 'a';
    }
}
\end{lstlisting}
\footnotesize
\begin{itemize}\itemsep=2pt
    \item interfaccia: \texttt{QUANTE}, \texttt{lettera}, \texttt{conta}. \texttt{isLettera} e \texttt{indice} al client non servono: \texttt{private}
    \item \texttt{c - 'a'} è già un \texttt{int}; \texttt{'a' + indice} anche, e per farne un \texttt{char} serve il cast
    \item perché non \texttt{Character.isLetter}? Per una \texttt{e} accentata risponde \texttt{true}, e l'indice sarebbe $232 - 97 = 135$: fuori dall'array
    \item niente \texttt{Scanner} né stampe: il modulo conta e basta
\end{itemize}
\end{columns}
\end{frame}

% ========================================================================
\begin{frame}[fragile]{Una possibile soluzione (2/2) -- il programma \texttt{Frequenze.java}}
\begin{columns}[T]
\column{0.6\textwidth}
\begin{lstlisting}[basicstyle=\ttfamily\scriptsize]
    public static void main(String[] args) {
        Scanner in = new Scanner(System.in);
        System.out.print("Testo: ");
        String testo = in.nextLine();
        in.close();

        int[] conteggi = Lettere.conta(testo);
        int totale = Statistiche.somma(conteggi);
        int max = Statistiche.massimo(conteggi);

        for (int i = 0; i < Lettere.QUANTE; i++) {
            int n = conteggi[i];
            if (n > 0) {
                double perc = (double) n / totale * 100;
                System.out.printf("%c %3d %5.1f%%",
                        Lettere.lettera(i), n, perc);
                if (n == max) {
                    System.out.print("  <-- max");
                }
                System.out.println();
            }
        }
    }
\end{lstlisting}

\column{0.4\textwidth}
\begin{lstlisting}[style=shell,basicstyle=\ttfamily\scriptsize\color{white}]
$ javac Frequenze.java
$ java Frequenze
Testo: Java e la JVM
a   3  30.0%  <-- max
e   1  10.0%
j   2  20.0%
l   1  10.0%
m   1  10.0%
v   2  20.0%
\end{lstlisting}
\small
Tre righe al centro del \texttt{main} dicono \emph{cosa} fa il programma; il \emph{come} sta nei moduli.

\vspace{0.15cm}
\texttt{Statistiche} è nato per i voti e qui lavora sui conteggi delle lettere, senza una modifica: è il \textbf{riuso}.

\vspace{0.15cm}
Stesso output di \texttt{FrequenzeMonolite}: è cambiata la struttura, non il comportamento.
\end{columns}
\end{frame}

% ========================================================================
\section{Wrap-up}

\begin{frame}{Riepilogo: cosa abbiamo trattato}
\footnotesize
\begin{center}
\renewcommand{\arraystretch}{1.25}
\begin{tabular}{|l|l|}
\hline
\textbf{Concetto (unità di riferimento)} & \textbf{Dove compare oggi} \\
\hline
Conversioni implicite, promozione (U11) & \texttt{Widening}; \texttt{return o[meta]} in \texttt{mediana} \\
\hline
Cast esplicito, troncamento (U11) & \texttt{Narrowing}; \texttt{media}; le percentuali in \texttt{Frequenze} \\
\hline
\texttt{char} e Unicode, aritmetica dei caratteri (U11) & \texttt{Caratteri}; \texttt{indice} e \texttt{lettera} in \texttt{Lettere} \\
\hline
Classi wrapper, costanti e metodi \texttt{static} (U11) & \texttt{Wrapper} \\
\hline
Autoboxing, unboxing, \texttt{==} ed \texttt{equals} (U11) & \texttt{Boxing} \\
\hline
La classe \texttt{Character} (U11) & \texttt{InfoCarattere}; \texttt{toLowerCase} in \texttt{conta} \\
\hline
\texttt{Arrays}: \texttt{toString}, \texttt{copyOf}, \texttt{equals}, \texttt{fill} (U11) & \texttt{UsaArrays}; \texttt{ordinata} in \texttt{Statistiche} \\
\hline
\texttt{Arrays}: \texttt{sort}, \texttt{binarySearch} (U11) & \texttt{OrdinaCerca}; \texttt{mediana} \\
\hline
Argomenti a lunghezza variabile (U11) & \texttt{Varargs}; tutti i metodi di \texttt{Statistiche} \\
\hline
Astrazione, modulo, interfaccia e corpo (U3) & \texttt{Voti} e \texttt{ProvaVoti} \\
\hline
Information hiding, \texttt{private} (U3) & \texttt{arrotonda}; \texttt{ordinata}; \texttt{isLettera} e \texttt{indice} \\
\hline
Coesione, accoppiamento, riuso (U3) & da \texttt{FrequenzeMonolite} a \texttt{Lettere} + \texttt{Statistiche} + \texttt{Frequenze} \\
\hline
\end{tabular}
\end{center}

\vspace{0.1cm}
\textbf{Prossima lezione (L5)}: i concetti della programmazione a oggetti e le \textbf{classi} in Java: incapsulamento, costruttori, metodi \texttt{get} e \texttt{set}, ciclo di vita degli oggetti.
\end{frame}

% ========================================================================
\begin{frame}{Homework 4 (da svolgere entro la prossima lezione)}
\begin{enumerate}\itemsep=5pt
    \item \textbf{\texttt{Anagrammi.java}}: legge due frasi e dice se sono una l'anagramma dell'altra, ignorando spazi, punteggiatura e maiuscole (\emph{``Calendario''} e \emph{``locandiera''} lo sono). Un metodo \texttt{char[] lettereOrdinate(String s)} \emph{crea} l'array delle sole lettere, rese minuscole con \texttt{Character} e ordinate con \texttt{Arrays.sort}; il confronto si fa con \texttt{Arrays.equals}.
    \item \textbf{\texttt{Pagella.java}}: i voti arrivano dalla riga di comando (\texttt{java Pagella 28 30 24 27 30 18}). Un metodo \texttt{int[] inInteri(String[] testi)} li converte con \texttt{Integer.parseInt}; il programma stampa i voti, i voti ordinati (su una copia), la media con 2 decimali, la mediana, la media arrotondata all'intero e la base di laurea in 110-esimi. Usate il modulo \texttt{Statistiche} di oggi \textbf{senza modificarlo}: copiatelo nella cartella \texttt{homework/}.
    \item \textbf{\texttt{Cifrario.java}} e \textbf{\texttt{ProvaCifrario.java}}: il cifrario di Cesare come modulo, senza la stringa \texttt{ALFABETO}: solo aritmetica dei \texttt{char} e cast. Interfaccia: \texttt{cifra(String, int)} e \texttt{decifra(String, int)}; nel corpo un metodo \texttt{private} \texttt{sposta(char, int)} che conserva maiuscole e minuscole. Il programma di prova legge testo e chiave, stampa il testo cifrato e quello decifrato, poi le 26 decifrature possibili del testo cifrato (\emph{forza bruta}).
\end{enumerate}

\vspace{0.2cm}
\small Per ogni metodo decidete se è interfaccia (\texttt{public}) o corpo (\texttt{private}). Usate solo ciò che abbiamo visto finora.
\end{frame}

% ========================================================================
\begin{frame}
\begin{center}
{\Huge Domande?}\\[2.5em]
{\normalsize \href{mailto:mattia.fonisto@unina.it}{\texttt{mattia.fonisto@unina.it}}}
\end{center}
\end{frame}

\end{document}
